\documentclass[../main.tex]{subfiles}
\begin{document}
%==========================================TIÊU ĐỀ TẬP========================================
\begin{table}[h]
  \begin{adjustwidth}{-1cm}{}
    \begin{tabular}{>{\centering\arraybackslash}p{6cm}|>{\raggedright\arraybackslash}p{12.5cm}}
      \multirow{5}{6cm}
      {\\[-40pt]
        \flaregothic\fontsize{20pt}{20pt}\selectfont \centering
        \color{eptype}HISTOIRE\color{black}\\
        \freshbot\fontsize{72pt}{72pt}\selectfont
        \color{epnum}42\\[6pt]
        \cafeteria\fontsize{20pt}{20pt}\selectfont\color{epnum}(D-42)\color{black}
        \\[18pt]
        \color{black}
      }
       &
      % Dòng thứ nhất: tên tập tiếng Nhật
      {\fotskip\fontsize{18pt}{18pt}\selectfont \ruby{苛}{いら}\ruby{立}{だ}ちの\ruby{裏}{うら}\ruby{側}{がわ}}
      \\[6pt]
      % Dòng thứ hai: tên tập tiếng Anh
       & {\cafeteria\fontsize{18pt}{18pt}\selectfont Behind the Irritation}                                                                                                     \\[4pt]
      % Dòng thứ ba: tên tập tiếng Việt
       & {\vnmsans\fontsize{18pt}{18pt}\selectfont Đằng sau sự cáu gắt}                                                                                                         \\[8pt]
      % Dòng thứ tư: Nhân vật xuất hiện
       & {\caviar{\fontsize{13pt}{13pt}\selectfont Nhân vật: \emp{kuon2} \emp{kaiho2} \emp{ryuuhou2} \emp{suzuno2} \emp{naomi2} \emp{game} \emp{roberu} \emp{oga} \emp{temma}}}
      \\[8pt]
    \end{tabular}
  \end{adjustwidth}
\end{table}
%=========================================TRUYỆN CHÍNH========================================
\color{black}

\txtbx[2]{
  {\color{vol} \uline{Tóm tắt:}} Những ngày sau vụ trộm mộ kinh hoàng và sự biến mất không dấu vết của Tô cùng Ki đẩy gia đình Hikawa vào một bầu không khí u uất và ngột ngạt đến nghẹt thở. Giữa lúc bố trầm mặc bên góc xưởng vắng, Suzuno vùi đầu tự trách trong phòng lab, Ryuuhou trút giận vào những bản rock và tờ rơi tìm cún vô vọng, thì Kaiho - người chị lớn kiêm Nữ tổng quản - lại trở thành tâm điểm của sự căng thẳng. Gánh nặng trách nhiệm quá lớn cùng nỗi bất lực trước kẻ thù vô hình khiến thần kinh của cô lúc nào cũng căng như dây đàn, biến Kaiho thành một ngọn núi lửa sẵn sàng bốc hỏa trước bất kỳ sự cố nhỏ nhặt nào. Từ những trận lôi đình trút xuống đầu các chàng trai Holostars bên mancave cho tới những giọt nước mắt kìm nén trong đêm muộn, Kuon từng bước thấu cảm nỗi đau sâu kín đằng sau lớp vỏ bọc gai góc của cô em gái song sinh.
}

\pov{Hikawa Kuon}

Có những nỗi đau không ồn ào gào thét, mà âm thầm gặm nhấm tâm can người ta qua từng ngày trôi trong sự bế tắc đến nghẹt thở.

Một tuần lễ đã trôi qua kể từ buổi sáng định mệnh khi nấm mồ của Mi bị đào xới tung tóe và hai cô cún nhỏ Tô, Ki hoàn toàn bốc hơi khỏi khu xưởng ngoại ô.

Khoảng thời gian ấy đối với cả gia đình tôi dài tựa như cả một thế kỷ. Những đợt tìm kiếm rầm rộ ban đầu trên khắp bốn phường của thủ đô Kawachi dần dần rơi vào ngõ cụt tăm tối. Hàng trăm tờ rơi in hình hai chú cún lông trắng dán trên các cột điện, trạm xe buýt và các bức tường gạch quanh khu phố Anwa bắt đầu bị cái nắng gắt mùa hạ làm cho cong queo, bạc màu. Những cuộc gọi báo tin từ người đi đường cứ thưa thớt dần, và cuối cùng thì tắt ngấm hoàn toàn.

Khu ký túc xá và xưởng cơ khí của gia đình chìm vào một bầu không khí nặng nề, ảm đạm chưa từng thấy.

Bố tôi trở nên trầm lặng đến đáng sợ. Mỗi buổi chiều tan ca ở xưởng cơ khí, người đàn ông dạn dày sương gió ấy lại lặng lẽ ra ngồi bệt trên bậc thềm hiên, đăm đăm nhìn về phía tảng đá hoa cương trắng lẻ loi dưới gốc cây rẻ quạt. Điếu thuốc lá trên môi ông cứ cháy đỏ lập lòe trong bóng chiều tà, làn khói đục ngầu bay lên che đi đôi mắt đục ngầu nỗi xót xa. Sau một cuộc họp gia đình ngắn ngủi trong nước mắt, bố đã đưa ra một quyết định dứt khoát và đau đớn: gia đình Hikawa sẽ vĩnh viễn không bao giờ nuôi thêm bất kỳ một con chó nào nữa. Nỗi đau mất đi Mi sau mười bốn năm gắn bó đã đủ xé lòng, nay vết thương ấy lại bị khoét sâu bởi sự mất tích bí ẩn của Tô và Ki, chẳng một ai trong nhà còn đủ dũng khí để đối diện với sự chia ly thêm một lần nào nữa.

Suzuno là người chịu tổn thương tâm lý nặng nề nhất. Cô thiếu nữ mười sáu tuổi tự giam mình trong gian phòng thí nghiệm mới dựng ở bãi đất hoang, vùi đầu vào các phản ứng hóa sinh phức tạp và các bản vẽ sơ đồ hệ thống cảm biến quang học an toàn. Em ấy tự dằn vặt bản thân mình suốt ngày đêm, cho rằng nếu chiều hôm ấy em không để Tô và Ki ngủ lại trông lab, hoặc nếu em cẩn thận hoàn thiện hệ thống an ninh sớm hơn một bước, thì thảm kịch đã không xảy ra. Mỗi lần dọn dẹp gian bếp tầng năm, nhìn thấy hai chiếc bát ăn bằng inox của cún vẫn nằm chỏng chơ bên mép tủ, khóe mắt dị sắc của Suzuno lại đỏ hoe, đôi bờ vai mảnh khảnh khẽ run lên trong câm lặng.

Ryuuhou thì phản ứng bằng sự bức bối và bướng bỉnh quen thuộc. Con bé dồn toàn bộ sự u uất vào dàn trống và những giai điệu rock giằng xé thâu đêm cùng ban nhạc underground. Mỗi buổi chiều tan học, Ryuuhou lại lủi thủi đạp xe đi dán lại từng tờ rơi bị rách, sục sạo khắp các bãi phế liệu và các trạm cứu hộ động vật đến mức khản đặc cả giọng. Sự ức chế tích tụ khiến con bé trở nên cực kỳ nhạy cảm và toxic mỗi khi chơi game hay lướt mạng xã hội, đến mức tôi phải liên tục đóng vai một chiếc van xả, ngồi bên cạnh đút kẹo ngọt và vỗ về để ghìm lại cơn bốc hỏa trong lòng em gái.

Bé út Naomi thỉnh thoảng lại lặng lẽ ôm chiếc vòng cổ cũ bằng vải mềm của Tô, ngước đôi mắt to tròn ngây thơ nhìn tôi hỏi: ``Anh Kuon ơi\dots\ bao giờ Tô với Ki mới nhớ đường về nhà ăn cơm vậy anh\dots?''

Mỗi câu hỏi ngây ngô ấy cất lên lại như một mũi kim đâm nhói vào lồng ngực tôi.

Thế nhưng, người chịu áp lực khủng khiếp nhất và có biểu hiện tâm lý biến chuyển nặng nề nhất trong ngôi nhà này\dots\ lại không phải là Suzuno hay Ryuuhou.

Mà đó chính là cô em gái song sinh của tôi: Hikawa Kaiho.

Khác với các em gái có thể khóc lóc, cáu kỉnh hay dồn nén vào âm nhạc để giải tỏa cảm xúc, Kaiho với tư cách là người chị lớn và là Nữ tổng quản điều hành của cả đại gia đình buộc phải tự khoác lên mình một lớp vỏ bọc sắt đá.

Cô vừa đau đớn, căm phẫn trước vụ trộm mộ báng bổ, vừa lo lắng thắt ruột cho tình trạng sức khỏe và tinh thần của Suzuno, vừa phải căng mình quán xuyến toàn bộ sổ sách tài chính, chứng từ kiểm kê của gia đình lẫn xưởng cơ khí ngoại ô.

Và trên hết, điều dày vò Kaiho nhiều nhất chính là sự bất lực của sức mạnh thể chất.

Bản thân cô sở hữu một sức mạnh cơ bắp phi thường đến mức có thể đấm nát một tảng đá tảng hay nhấc bổng cả một chiếc xe tải nhẹ. Thế nhưng, trước một kẻ thù vô hình thoắt ẩn thoắt hiện, một vụ án không để lại bất kỳ một dấu vết hay mẫu sinh trắc học nào trên toàn bộ mạng lưới dữ liệu của Game, nắm đấm thép của Kaiho hoàn toàn rơi vào khoảng không vô định. Cô không biết phải tìm ai để đòi lại công lý, không biết phải trút cơn thịnh nộ đang sôi sục trong huyết quản vào đâu.

Sự bất lực và quá tải tinh thần ấy đã biến Kaiho thành một ngọn núi lửa âm ỉ chực chờ phun trào bất kỳ lúc nào.

Dạo gần đây, thần kinh của cô lúc nào cũng căng như một sợi dây đàn sắp đứt. Đôi mày thanh tú của Kaiho luôn nhíu chặt lại thành một rãnh sâu, ánh mắt xanh thẫm lúc nào cũng đằng đằng sát khí và cực kỳ nhạy cảm trước mọi biến động xung quanh.

Ở ký túc xá phố Anwa, bầu không khí luôn trong tình trạng báo động cấp cao. Ai vô tình đóng cánh cửa hơi mạnh tay một chút, ai để chiếc dép lê lệch khỏi vạch sàn, hay chỉ cần một giọt nước vô ý rớt ra mặt bàn chưa kịp lau khô\dots\ đều phải hứng chịu một cái lườm sắc như dao cạo và những lời răn đe lạnh tanh đến thấu xương từ nàng Tổng quản. Cả nhà ai nấy đều phải đi nhẹ, nói khẽ, cố gắng không chạm vào ngòi nổ nguy hiểm đang chực bốc cháy kia.

Thế nhưng, nơi phải nếm trải cơn thịnh nộ ấy nặng nề và dở khóc dở cười nhất\dots\ lại chính là căn mancave của các chàng trai \hls.

Là người phụ trách giám sát chi tiêu, an toàn cơ sở vật chất và lịch sinh hoạt của nhóm nam, Kaiho thường xuyên phải sang căn hầm mancave để kiểm tra định kỳ. Và trong những ngày tâm trạng đen tối này, căn mancave vốn ồn ào náo nhiệt bỗng chốc biến thành một thao trường kỷ luật sắt.

Tôi nhớ như in buổi chiều hôm trước ngày diễn ra sự việc kỳ quặc của Roberu. Khi tôi cùng Kaiho bước vào căn phòng hầm với xấp hồ sơ kiểm kê trên tay, bầu không khí đã lập tức chùng xuống.

Chàng Hiệp sĩ Temma đang vui vẻ pha một tách cà phê thơm lừng bên quầy bar nhỏ. Vì mải ngâm nga một điệu hát, cậu lóng ngóng để rớt đúng một giọt cà phê nâu sẫm xuống mặt bàn kính vừa được lau sạch bóng.

Ngay khoảnh khắc giọt nước chạm mặt kính phát ra một tiếng *TÁCH* siêu nhỏ, Kaiho đang ngồi bên bàn làm việc bỗng dằn mạnh cây bút bi xuống mặt bàn phát ra một tiếng *CỘP* đanh gọn.

Cô ngẩng phắt dậy, đôi mắt xanh lam phóng ra những tia nhìn sắc lẹm như muốn đâm thủng người đối diện. Temma giật bắn mình, chiếc thìa khuấy cà phê suýt nữa rơi khỏi tay, cậu chàng đứng chết trân tại chỗ, mặt mày tái mét.

Và thế là, suốt hai mươi phút đồng hồ sau đó, cả căn mancave chìm vào một sự tĩnh lặng chết chóc, chỉ còn lại giọng nói lanh lảnh, rành rọt và sắc như lưỡi dao của Kaiho vang vọng khắp các ngóc ngách. Cô bắt Temma phải đứng nghiêm hai tay khoanh trước ngực, giảng giải cặn kẽ từ nguyên tắc bảo dưỡng mặt kính cường lực, tác hại của axit hữu cơ trong cà phê đối với lớp phủ nano, cho đến tinh thần trách nhiệm và tính cẩn trọng tối thiểu của một người trưởng thành.

Temma đứng nghe mà mồ hôi hột túa ra như tắm, hai chân run lẩy bẩy, chỉ biết liên tục gật đầu lia lịa: ``Vâng\dots\ Em xin lỗi Kaiho-senpai\dots\ Lần sau em tuyệt đối không để rớt một giọt nào nữa đâu ạ\dots''

Chưa dừng lại ở đó, đến lượt Oga bước vào nộp báo cáo. Do bận rộn với lịch huấn luyện thể lực dày đặc vào buổi sáng, Oga nộp chậm hóa đơn tiền điện và bảng kê bảo trì phòng tập chỉ vỏn vẹn nửa buổi.

Kaiho liếc nhìn tờ hóa đơn, rồi ngước mắt nhìn lên thân hình hộ pháp cơ bắp cuồn cuộn của Oga với một nụ cười lạnh tanh khiến sống lưng gã quân nhân cũng phải ớn lạnh: ``Trễ ba tiếng bốn mươi hai phút. Cậu nghĩ hệ thống lưới điện công cộng của thủ đô này sẽ ngồi chờ các cậu tập cơ bắp xong rồi mới đến ngày thu tiền chắc? Cậu có ba mươi phút để hoàn tất toàn bộ các chữ ký xác nhận còn thiếu. Bằng không, đích thân tôi sẽ dập cầu dao tổng của phòng tập, cho toàn bộ các cậu rèn luyện thể lực trong bóng tối giữa trưa hè ba mươi tám độ C này.''

Oga nuốt nước bọt cái ực, khuôn mặt nghiêm nghị thường ngày lộ rõ vẻ hoang mang tột độ, cậu vội vã ôm lấy tập hồ sơ lao đi ký tên như một người lính vừa nhận lệnh hành quân khẩn cấp.

Phía sau thành ghế sofa, Rikka, Astel và Roberu nấp thụp người xuống, chỉ dám thò nửa con mắt ra quan sát, thì thào với nhau trong sự khiếp đảm:

``Trời đất ơi\dots\ Dạo gần đây Kaiho-senpai bị làm sao thế nhỉ? Sao chị ấy dữ dằn như chằn tinh thức giấc vậy trời\dots''

``Tuần trước em lỡ để đôi giày tập lệch một góc mười độ thôi mà cũng bị chị ấy lườm cho xém rách cả áo khoác đấy\dots''

``Ai ở nhà chọc giận bà chằn lửa này thế không biết\dots\ Khéo thở mạnh một cái thôi cũng bị chị ấy táng cho bay màu mất thôi\dots''

Tôi đứng tựa lưng bên khung cửa, nhìn cảnh tượng dở khóc dở cười ấy mà chỉ biết lắc đầu thở dài ngao ngán, kín đáo ra hiệu cho đám đàn em hãy ngoan ngoãn ngậm miệng lại để bảo toàn tính mạng.

Tôi hiểu rõ hơn ai hết: Kaiho không hề ác ý, cũng chẳng phải là người thích bắt nạt kẻ khác. Đó chỉ là phản ứng tự vệ trong vô thức của một tâm hồn đang chịu đựng quá nhiều thương tổn và gánh nặng mà không tìm được lối thoát.

\ct

Trở về ký túc xá phố Anwa khi ráng chiều mùa hạ đã tắt hẳn.

Sau một ngày dài mỏi mệt bên xưởng cơ khí và căn mancave, tôi bước lên các tầng lầu của khu nhà, nơi những đứa em gái của tôi cũng đang âm thầm trải qua những cuộc chiến nội tâm giằng xé của riêng mình.

Điểm dừng chân đầu tiên của tôi là căn phòng riêng của Ryuuhou.

Vừa bước tới đầu hành lang, những âm thanh giằng xé, chát chúa của tiếng đàn guitar điện kết hợp với nhịp trống dồn dập vang ra từ khe cửa phòng con bé đã đập vào tai tôi. Tiếng nhạc rock không mang vẻ ngẫu hứng của một buổi tập thông thường, mà chứa chan sự giận dữ, bất lực và u uất bị kìm nén đến nghẹt thở.

Tôi khẽ vặn nắm đấm cửa bước vào.

Ryuuhou đang ngồi bệt dưới sàn gỗ, dựa lưng vào mép giường. Cây guitar điện để nằm nghiêng bên cạnh, trên màn hình máy tính vẫn còn hiển thị dòng chữ đỏ rực báo hiệu một ván đấu game thất bại thảm hại. Đôi giày thể thao trắng lấm lem bùn đất bị con bé cởi vứt bừa bãi nơi góc phòng, xấp tờ rơi tìm cún nhăn nhúm nằm rải rác trên bàn học.

Cô thiếu nữ mười sáu tuổi ngồi ôm hai đầu gối, mái tóc ngắn rối bời, đôi mắt đỏ ngầu đăm đăm nhìn trừng trừng vào khoảng không trước mặt.

Tôi không nói gì, lẳng lặng bước tới nhặt xấp tờ rơi xếp lại ngay ngắn trên bàn, rồi rút từ trong túi áo ra một phong kẹo ngậm vị bạc hà mát lạnh và một chai nước khoáng, ngồi bệt xuống sàn gỗ cạnh em gái.

``Uống chút nước đi,'' tôi đưa chai nước cho con bé.

Ryuuhou giật mình, liếc nhìn tôi rồi buông một tiếng thở hắt đầy vẻ bực bội: ``Anh Kuon, anh vào đây làm gì? Em đang bực mình, đừng có làm phiền em.''

``Biế là em bực rồi,'' tôi ôn tồn mở nắp chai nước đặt vào tay em. ``Hôm nay lại đạp xe đi khắp các bãi phế liệu bên quận Nam đúng không?''

Nghe tôi nhắc đến chuyện tìm cún, sống mũi Ryuuhou bỗng cay xè. Con bé siết chặt nắm tay đập mạnh xuống sàn gỗ một tiếng *BỊCH* nghẹn ức: ``Chẳng có tin tức gì cả\dots! Chỗ nào người ta cũng xua đuổi, bảo là chó cỏ mất tích thì tìm làm cái gì cho mất công! Cột điện nào em dán tờ rơi thì chỉ được nửa ngày là bị người ta xé toạc hoặc bị nắng mưa làm cho nát bét\dots! Khốn kiếp thật chứ, rốt cuộc là kẻ nào đã làm chuyện này\dots\ Sao bọn chúng lại nhẫn tâm đến mức đào xới cả mộ của con Mi lên cơ chứ!?''

Hai hàm răng con bé bặm chặt bờ môi đến rớm máu, những giọt nước mắt uất ức thi nhau lăn dài trên gò má lấm lem bụi đường.

Tôi đưa tay xoa nhẹ mái tóc ngắn bù xù của em gái, kéo con bé tựa nhẹ vào vai mình: ``Em đã làm hết sức mình rồi, Ryuuhou. Mi, Tô và Ki nếu nhìn thấy em vất vả vì chúng nó như thế, chắc chắn cũng sẽ rất xót xa đấy.''

Tôi bóc một viên kẹo ngọt đút vào miệng em gái, cất giọng dịu dàng nhưng đầy kiên định: ``Nhưng đừng biến nỗi đau thành cơn giận dữ mù quáng rồi tự làm tổn thương chính mình. Anh hiểu em muốn gào thét, muốn đập phá cho hả cơn tức, nhưng nếu em cứ tự hành hạ bản thân bằng những đêm thức trắng và những ván game cay cú như thế này, ai sẽ cùng anh tiếp tục đi tìm Tô và Ki đây?''

Ryuuhou ngậm viên kẹo ngọt, vị the mát lan tỏa nơi đầu lưỡi làm dịu bớt ngọn lửa đang sôi sục trong lồng ngực. Con bé khẽ sụt sịt mũi, gục đầu vào vai tôi, hai tay níu chặt lấy vạt áo anh trai: ``Em nhớ hai đứa nó lắm anh ơi\dots\ Chiều nào đi học về, hai đứa nó cũng chạy xổ ra cổng đón em cơ mà\dots''

``Anh biết, anh cũng nhớ chúng nó,'' tôi vỗ nhẹ lưng em gái. ``Ngoan ngoãn đi tắm rửa rồi nghỉ ngơi sớm đi. Anh hứa với em, chỉ cần còn một tia hy vọng, anh nhất định sẽ không bao giờ bỏ cuộc.''

\ct

Rời khỏi phòng Ryuuhou, tôi bước sang góc hành lang đối diện, nơi cánh cửa phòng thí nghiệm nhỏ kiêm phòng riêng của Suzuno vẫn khép hờ.

Ánh đèn neon trắng lạnh lẽo hắt qua khe cửa.

Suzuno đang ngồi cô độc bên chiếc bàn vi lượng, đôi mắt dán chặt vào thị kính của chiếc kính hiển vi quang học. Trên mặt bàn chi chít những bản vẽ phác thảo sơ đồ cảm biến laser đa điểm và các mẫu dung dịch đang phân tích dở dang. Chiếc khay đựng phần cơm tối mà mẹ tôi mang lên từ hai tiếng trước vẫn nằm nguyên vẹn bên góc bàn, đĩa thức ăn đã nguội ngắt tự bao giờ.

Tôi khẽ gõ cửa bước vào, mang theo một cốc sữa ấm pha chút mật ong thơm phức.

``Suzuno à, nghỉ tay một chút đi em.''

Suzuno giật mình ngẩng đầu lên. Gương mặt thanh tú của cô bé mười sáu tuổi trông xanh xao và gầy rộc đi trông thấy sau những ngày thức trắng. Nhìn thấy tôi, cô bé vội vã đưa tay gạt những lọn tóc lòa xòa trước trán, ánh mắt thoáng lộ vẻ bối rối tội lỗi: ``Anh Kuon\dots\ Em xin lỗi, em đang kiểm tra lại cấu trúc phân tử của vệt dung môi hữu cơ thu thập được ở hiện trường nấm mồ\dots\ Em muốn tìm ra công thức chính xác của chất bảo quản mô mà kẻ trộm mộ đã sử dụng\dots''

Tôi nhẹ nhàng đặt cốc sữa ấm xuống bàn, đẩy chiếc khay cơm nguội sang một bên rồi kéo chiếc ghế xoay ngồi xuống cạnh em gái: ``Chuyện đó thì để ngày mai làm tiếp. Cả ngày nay em đã ăn gì vào bụng chưa?''

Suzuno cúi gằm mặt, hai bàn tay nhỏ nhắn bấu chặt vào tà áo blouse trắng, giọng nói nghẹn lại như sắp vỡ vụn: ``Em\dots\ em không đói\dots\ Em nuốt không trôi anh ơi\dots''

Một giọt nước mắt lặng lẽ rơi xuống tờ giấy ghi chép cấu trúc hóa học, làm nhòe đi những nét mực xanh: ``Là tại em cả\dots\ Giá như chiều hôm ấy em không để Tô và Ki ở lại xưởng trông lab\dots\ Giá như em cẩn thận lắp đặt hệ thống cảnh báo sớm hơn một tuần\dots\ thì mẹ Mi đã không bị kẻ xấu quấy rầy, Tô và Ki cũng đã không bị bắt đi mất\dots\ Em là một đứa tồi tệ, em đã không bảo vệ được hai đứa nhỏ\dots''

Nhìn những vết trầy xước do gai dại cào rách trên mu bàn tay nhỏ nhắn của Suzuno vẫn còn vương những vệt thuốc sát trùng màu đỏ, lòng tôi xót xa đến thắt lại.

Tôi nắm lấy bàn tay lạnh ngắt của em gái, giọng nói trầm ấm và bao dung: ``Suzuno, nhìn anh này.''

Cô bé ngước đôi mắt ướt đẫm nhìn tôi.

``Kẻ dám đào trộm mộ và bắt cóc Tô, Ki là những kẻ có tổ chức, có toan tính và sở hữu công nghệ tinh vi đến mức ngay cả hệ thống quét dữ liệu của Game cũng bị phong tỏa,'' tôi nghiêm túc phân tích cho em hiểu. ``Đó không phải là một vụ trộm vặt thông thường, và một cô bé mười sáu tuổi như em hoàn toàn không thể nào lường trước được điều đó. Việc em tự dằn vặt và bỏ đói bản thân không giúp ích được gì cho Tô và Ki, mà chỉ khiến cả nhà thêm đau lòng thôi.''

Tôi đưa cốc sữa ấm vào đôi bàn tay nhỏ bé của em: ``Em có biết không, chính cái mũi sắc ký sống và kiến thức hóa sinh của em mới là niềm hy vọng lớn nhất để chúng ta lật mở chân tướng vụ án này. Nếu em gục ngã vì kiệt sức, ai sẽ giúp anh phân tích những vệt dung môi bí ẩn kia? Uống hết cốc sữa này rồi ăn chút cơm đi, vì cả nhà, và vì chính Tô, Ki nữa, được không em?''

Nghe những lời phân tích thấu tình đạt lý và sự tin tưởng tuyệt đối của người anh cả, tảng đá đè nặng trong lòng Suzuno như vỡ vụn. Cô bé khẽ gật đầu, hai hàng nước mắt tuôn rơi nhưng ánh mắt đã vơi đi phần nào nỗi ám ảnh tội lỗi: ``Vâng\dots\ Em nghe lời anh\dots\ Em sẽ cố gắng ăn hết\dots''

\ct

Khoảng mười một giờ đêm.

Trước khi lên tầng năm, tôi ghé qua phòng ngủ của bé út Naomi ở tầng ba.

Cánh cửa phòng mở hé. Trong gian phòng mờ ảo ánh đèn ngủ hình chú cá voi, bé út vẫn chưa chịu đi ngủ. Con bé đang ngồi bó gối trên giường, hai cánh tay nhỏ xíu ôm chặt lấy chú gấu bông xù lông và chiếc vòng cổ bằng vải mềm màu đỏ của Tô. Đôi mắt to tròn trong veo của con bé đăm đăm nhìn ra khung cửa sổ mở hé, nơi gió đêm thổi qua rèm cửa khẽ đung đưa.

Tôi rón rén bước vào, ngồi xuống bên mép giường, khẽ vuốt ve mái tóc mềm mại của em gái: ``Naomi, muộn lắm rồi, sao em chưa chịu ngủ thế?''

Naomi ngước nhìn tôi, đôi mắt ngấn lệ chớp chớp, giọng thỏ thẻ non nớt khiến lồng ngực tôi thắt lại: ``Anh Kuon ơi\dots\ Tô với Ki có bị đói không anh? Ở ngoài trời đêm lạnh lắm, ai sẽ đắp chăn cho hai đứa nó ngủ vậy anh\dots? Tô sợ bóng tối lắm, mỗi lần sấm chớp là nó lại chui vào gầm giường run cầm cập cơ mà\dots''

Những câu hỏi ngây thơ, thuần khiết của đứa trẻ chín tuổi luôn chạm vào những góc đau đớn nhất trong tim người lớn.

Tôi nhẹ nhàng ôm lấy thân hình nhỏ nhắn của Naomi vào lòng, để đầu con bé tựa vào ngực mình, tay vỗ về nhịp nhàng sau lưng em: ``Tô và Ki rất thông minh và dũng cảm, Naomi à. Hai đứa nó được bố dạy dỗ từ nhỏ, lại được em và Suzuno chăm sóc cẩn thận, nhất định chúng nó sẽ biết cách tìm nơi trú ẩn an toàn và ấm áp nhất.''

Tôi nâng chiếc vòng cổ vải đỏ trong tay con bé lên, cẩn thận gấp lại đặt ngay ngắn trên đầu giường: ``Anh hứa với Naomi, dù có phải đi tìm khắp nơi, dù có mất bao lâu đi chăng nữa, anh và cả nhà nhất định sẽ đưa Tô và Ki trở về bên em. Nên bây giờ Naomi phải ngoan ngoãn ngủ một giấc thật ngon, để ngày mai còn có sức đón hai đứa nó về, em đồng ý không?''

Naomi ngước nhìn tôi, ánh mắt thơ ngây dần tìm lại được cảm giác an tâm. Con bé khẽ gật đầu: ``Dạ\dots\ Em tin anh Kuon\dots''

Tôi nhẹ nhàng kéo chăn đắp ngang ngực cho em gái, đặt chú gấu bông vào lòng con bé rồi khe khẽ ngâm nga một giai điệu ru êm dịu. Chỉ vài phút sau, hơi thở của Naomi đã trở nên đều đặn, con bé chìm vào giấc ngủ yên bình với đôi môi khẽ mỉm cười.

Tôi hôn nhẹ lên trán em gái, rồi rón rén khép hờ cánh cửa bước ra ngoài.

\ct

Sau khi đã đi một vòng lắng nghe, xoa dịu và chia sẻ nỗi đau với ba đứa em gái, tôi bước từng bước chậm rãi lên cầu thang dẫn lên tầng năm.

Đêm hôm đó, đồng hồ trên vách tường tầng năm điểm mười hai giờ khuya.

Cả khu ký túc xá đã chìm sâu vào giấc ngủ yên tĩnh sau một ngày dài mệt mỏi. Thế nhưng, qua khe cửa phòng làm việc chung ở cuối hành lang, ánh đèn bàn màu vàng nhạt vẫn hắt ra một vệt sáng lẻ loi.

Tôi khẽ đẩy cánh cửa gỗ bước vào, trên tay bưng một chiếc khay nhỏ với hai tách trà hoa cúc mật ong nóng hổi và một đĩa bánh quy bơ giòn rụm.

Ngồi bên chiếc bàn gỗ dài, Kaiho vẫn chưa chịu đi ngủ. Trước mặt cô là những chồng sổ sách dày cộp, những tập hóa đơn chứng từ kiểm kê của xưởng cơ khí, và nằm ngay bên cạnh góc bàn\dots\ là một xấp tờ rơi tìm kiếm Tô và Ki còn sót lại chưa kịp dán.

Cô em gái song sinh của tôi ngồi đó, hai khuỷu tay chống lên mặt bàn, mười đầu ngón tay luồn sâu vào mái tóc màu đen, siết chặt lấy hai bên thái dương. Bờ vai cô hơi khom xuống, để lộ sự rệu rã và kiệt quệ sau những chuỗi ngày gồng mình chống đỡ.

Tôi nhẹ nhàng bước tới, đặt tách trà bốc khói nghi ngút xuống chiếc lót ly bằng gỗ trước mặt cô: ``Uống chút trà nóng đi em. Muộn lắm rồi đấy.''

Kaiho giật mình khẽ ngẩng đầu lên. Nhìn thấy tôi, đôi mày cô hơi chau lại theo thói quen phòng thủ, giọng nói mang chút gắt gỏng mỏi mệt: ``Anh chưa ngủ à? Em đang bận đối chiếu lại các khoản chi phí sửa chữa máy móc ở xưởng, anh đừng có làm phiền.''

Tôi không đáp lời càu nhàu ấy. Tôi kéo chiếc ghế đối diện ngồi xuống, từ tốn nhấp một ngụm trà ấm, rồi nhìn thẳng vào đôi mắt màu xanh thẫm đang trũng sâu vì thiếu ngủ của em gái.

Ánh mắt ấy chứa chan sự giằng xé giữa vẻ đanh thép bên ngoài và sự yếu mềm đang rỉ máu bên trong.

Tôi cất giọng trầm ấm, nhẹ nhàng phá vỡ bức tường phòng ngự kiên cố: ``Em đang sợ hãi đúng không, Kaiho?''

Bàn tay đang cầm bút bi của Kaiho bỗng khựng lại giữa không trung. Cô ngước mắt lên nhìn tôi, ánh nhìn lóe lên một tia bực bội: ``Anh nói cái quái gì thế? Em sợ cái gì chứ?''

``Em sợ rằng bản thân mình không đủ năng lực để bảo vệ những người thân yêu trong ngôi nhà này,'' tôi điềm đạm tiếp lời, từng câu chữ găm thẳng vào góc khuất sâu kín nhất trong lòng em gái. ``Em cáu bẳn với bố mẹ, em nghiêm khắc gắt gỏng với Suzuno và Ryuuhou, em trút những trận lôi đình vô cớ xuống đầu mấy đứa nhỏ bên \hls\dots\ Tất cả những điều đó thực chất không phải vì em giận dữ với họ. Mà là vì em đang giận dữ và căm ghét chính bản thân mình.''

Kaiho cắn chặt bờ môi dưới, những ngón tay siết chặt lấy thân bút bi đến mức các khớp xương trắng bệch ra: ``Anh im đi\dots''

``Em giận vì bản thân sở hữu một sức mạnh thể chất áp đảo, có thể đấm nát cả một bức tường bê tông, nhưng lại hoàn toàn bất lực trước một chiếc hố đất trống rỗng dưới gốc cây rẻ quạt,'' tôi nhìn thẳng vào mắt cô, giọng nói chan chứa sự xót xa thấu hiểu. ``Em xót xa khi thấy Suzuno gầy rộc đi vì tự dằn vặt mình, em đau lòng khi nghe Naomi ngây thơ hỏi bao giờ cún mới về ăn cơm, và em lo sợ rằng kẻ trộm mộ kia là một thế lực nguy hiểm đang âm thầm rình rập gia đình này. Em gồng mình lên đóng vai một người bảo vệ thép, muốn một mình ôm hết mọi âu lo và công việc vào người, để rồi thần kinh bị quá tải đến mức sắp đứt tung ra.''

Ngòi bút bi trong tay Kaiho run rẩy kịch liệt, rồi rơi *CẠCH* một tiếng xuống mặt bàn gỗ.

Lớp vỏ bọc sắt đá và lạnh lùng mà cô cố gắng dựng lên suốt những ngày qua bỗng chốc sụp đổ hoàn toàn trước ánh mắt nhìn thấu tâm can của người anh trai song sinh.

Kaiho buông thõng hai tay xuống đùi, đầu cúi gằm, hai hàng nước mắt nóng hổi bất giác lăn dài trên gò má thanh tú. Giọng cô run lên bần bật, vỡ òa ra trong sự nghẹn ngào uất ức: ``Đúng thế đấy\dots! Em bất lực lắm anh Kuon ơi\dots! Nắm đấm này của em thì có tích sự gì chứ\dots? Sức mạnh quái vật này rốt cuộc để làm gì khi mà thi hài của con Mi bị người ta cướp mất em cũng không giữ được, hai đứa nhỏ Tô và Ki bị bắt cóc đi đâu em cũng không tìm ra\dots!''

Cô đưa hai bàn tay lên ôm chặt lấy mặt, những giọt nước mắt ấm nóng thấm ướt qua kẽ ngón tay: ``Nhìn Suzu cứ cặm cụi trong lab không chịu ăn uống, nhìn Nao cứ ôm cái vòng cổ cũ\dots\ lòng em đau như bị dao cứa vậy anh ơi\dots\ Em sợ lắm\dots\ Em sợ nếu kẻ thù nhắm vào gia đình mình, em sẽ lại chậm chân một bước và để mất thêm một ai đó nữa\dots''

Tôi im lặng lắng nghe tiếng khóc nghẹn ngào của em gái.

Đã bao lâu rồi tôi mới lại thấy Kaiho khóc như thế này? Đó là lần khi chúng tôi ở The Void, khi mà chỉ còn lại Suzuno và Kaiho, và cô em gái song sinh sau này của tôi đã tìm thấy một bến đỗ mới sau khi phải trốn chạy trong thế giới cũ đầy rẫy sự mạt sát của cánh đàn ông. Cô gái kiên cường này luôn giành phần gánh vác những nỗi đau để bảo vệ em mình, và sau này là cả gia đình tôi - một mái ấm mà cô đã hằng mong ước suốt tám năm.

Cô luôn tỏ ra mạnh mẽ, bất khả chiến bại, nhưng suy cho cùng, cô vẫn chỉ là một thiếu nữ mang trái tim tràn ngập tình yêu thương gia đình.

Tôi đứng dậy, bước vòng qua chiếc bàn gỗ, đưa tay vỗ nhẹ lên bờ vai đang run rẩy của em gái rồi kéo đầu cô tựa nhẹ vào người mình như những ngày thơ ấu.

``Ngốc ạ,'' tôi mỉm cười dịu dàng, khẽ xoa đầu cô. ``Hai chúng ta sinh ra cùng một ngày, cùng chia sẻ một dòng máu, là anh em song sinh cơ mà. Ai cho phép em tự tiện gánh hết toàn bộ gánh nặng của ngôi nhà này lên đôi vai một mình thế?''

Tôi nhìn xuống xấp hóa đơn và những tờ rơi tìm chó trên bàn: ``Chúng ta không biết kẻ đứng sau là ai, nhưng gia đình Hikawa chưa bao giờ đơn độc. Có anh ở đây, có bố mẹ, có các em, và có cả những người bạn luôn sẵn sàng sát cánh cùng chúng ta. Nắm đấm của em không hề vô dụng, nó là chiếc khiên vững chắc nhất bảo vệ ngôi nhà này. Nhưng em phải giữ cho cái đầu mình thật tỉnh táo, đừng để sự giận dữ thiêu rụi chính bản thân mình.''

Từ mặt chiếc đồng hồ thông minh trên cổ tay tôi, một đốm sáng xanh lam khẽ nhấp nháy êm ả. Giọng nói kim loại phẳng lặng của Game cất lên qua loa nhỏ, lần này không hề có sự mỉa mai thường ngày: ``Chỉ số cortisol và nồng độ adrenaline trong máu của Kaiho-user đang hạ dần về ngưỡng an toàn. Phân tích dữ liệu hành vi cho thấy: việc chia sẻ áp lực cảm xúc giúp gia tăng 42\% độ minh mẫn trong các quyết định quản trị tiếp theo. Tôi khuyến nghị hai người nên uống hết tách trà trước khi nhiệt độ chất lỏng giảm xuống dưới 40 độ C.''

Nghe thấy lời nhắc nhở đậm chất máy móc của vị Trợ lý ảo, Kaiho khẽ sụt sịt mũi, đưa mu bàn tay quẹt vội những giọt nước mắt trên má rồi bật cười trong sự ngượng ngùng: ``Đến cả ông cũng hùa vào dạy đời tôi đấy à, Game?''

Game đáp gọn lỏn: ``Tôi chỉ đang thực hiện chức năng chăm sóc sức khỏe cho người dùng theo đúng giao thức lập trình.''

Tôi khẽ cười, đưa tách trà hoa cúc còn ấm áp vào tay em gái: ``Ngoan ngoãn uống hết đi rồi đi ngủ sớm. Ngày mai anh sẽ sang mancave phụ em rà soát lại đống sổ sách bên đó, tiện thể nói chuyện để mấy đứa nhỏ bên ấy hiểu và thông cảm cho em hơn.''

Kaiho đón lấy tách trà, hai bàn tay áp chặt vào thành gốm ấm áp, hít một hơi thật sâu làn hương thảo mộc thanh nhẹ. Ánh mắt xanh thẫm của cô đã lấy lại được sự trong trẻo và kiên định vốn có.

Cô khẽ gật đầu, nở một nụ cười nhẹ nhõm: ``Vâng. Cảm ơn anh, anh Kuon.''

Màn đêm buông xuống bên ngoài khung cửa sổ phố Anwa vẫn lạnh lẽo và tĩnh mịch, nhưng bên trong căn phòng làm việc nhỏ bé, hơi ấm của tình thân đã xua tan đi sự ngột ngạt tích tụ suốt bao ngày qua.

Tôi biết rằng những thử thách cam go vẫn còn đang chờ đợi ở phía trước, và hành trình tìm lại những người bạn bốn chân vẫn còn mịt mờ trong giông bão. Thế nhưng, chỉ cần những sợi dây gắn kết của gia đình này vẫn được thắt chặt bằng sự thấu hiểu và sẻ chia chân thành, thì không một thế lực hắc ám nào có thể đánh gục được chúng tôi.

\end{document}